% GENERATED FILE. Edit canonical lessons, then run npm run generate:books.
% canonical-source-hash: fnv1a64:3dd5c3a29d3a0efa
% canonical-lessons: BN-C106-chhora, BN-C106-tepa, BN-C106-aka, BN-C106-phera, BN-C106-chalano

\chapter{To throw, To press, To draw, To return, To drive}
\label{ch:bn-to-throw-to-press-to-draw-to-return-to-drive}
\hlchaptermodality{106}

\begin{chapteropening}
\textbf{By the end of this chapter:} \emph{I can say to throw, to press, to draw, to return and to drive.}

Complete the last lesson of chapter 106: I can say to throw, to press, to draw, to return and to drive.
\end{chapteropening}

\section[\texorpdfstring{\bn{ছোড়া} (chhõṛā)}{ছোড়া (chhõṛā)}]{\bn{ছোড়া} (chhõṛā) — to throw}

\label{lesson:BN-C106-chhora}



\begin{quote}
Before the new one: say the Bengali for to fry, then the Bengali for to mix.
\end{quote}

\subsection*{You'll want to know: \bn{ছোড়া}}
\textbf{\bn{ছোড়া}} — \emph{chhõṛā} — \textquotedblleft{}to throw\textquotedblright{}.

\begin{cousinweb}
\textbf{\bn{ছোড়া}} (\emph{chhõṛā}) is a tadbhava, a word inherited from Sanskrit and reshaped by sound change. It comes from \emph{kṣoṭayati}, \textquotedblleft{}throws\textquotedblright{} (CDIAL 3747).
\end{cousinweb}

\begin{grammarlens}[title={What you've built}]
One more word for this chapter.
\end{grammarlens}

\subsection*{Your turn}
\begin{itemize}
\raggedright
  \item \emph{Say it:} \emph{chhõṛā}
  \item \emph{Say it:} \emph{chhõṛā}, once more
  \item \emph{Say it:} say \emph{chhõṛā}
  \item \emph{Recall:} say the Bengali for to fry, then the Bengali for to mix, then say \emph{chhõṛā} again
\end{itemize}

\begin{tcolorbox}[breakable,colback=teal!4,colframe=teal!35!black,title={Before you move on}]
What does \bn{ছোড়া} mean? (\textquotedblleft{}To throw\textquotedblright{}.) Say it once more.
\end{tcolorbox}
\section[\texorpdfstring{\bn{টেপা} (ṭepā)}{টেপা (ṭepā)}]{\bn{টেপা} (ṭepā) — to press}

\label{lesson:BN-C106-tepa}



\begin{quote}
Before the new one: say the Bengali for to mix, then the Bengali for to throw.
\end{quote}

\subsection*{You'll want to know: \bn{টেপা}}
\textbf{\bn{টেপা}} — \emph{ṭepā} — \textquotedblleft{}to press\textquotedblright{}.

\begin{cousinweb}
The origin of \textbf{\bn{টেপা}} (\emph{ṭepā}) is not settled. CDIAL has no entry for it; the nearest form it gives is Bengali \emph{ṭipā}, \textquotedblleft{}to pinch\textquotedblright{}, under \emph{*ṭipp} (CDIAL 5464), and it does not join the two.
\end{cousinweb}

\begin{grammarlens}[title={What you've built}]
One more word for this chapter.
\end{grammarlens}

\subsection*{Your turn}
\begin{itemize}
\raggedright
  \item \emph{Say it:} \emph{ṭepā}
  \item \emph{Say it:} \emph{ṭepā}, once more
  \item \emph{Say it:} say \emph{ṭepā}
  \item \emph{Recall:} say the Bengali for to mix, then the Bengali for to throw, then say \emph{ṭepā} again
\end{itemize}

\begin{tcolorbox}[breakable,colback=teal!4,colframe=teal!35!black,title={Before you move on}]
What does \bn{টেপা} mean? (\textquotedblleft{}To press\textquotedblright{}.) Say it once more.
\end{tcolorbox}
\section[\texorpdfstring{\bn{আঁকা} (ā̃kā)}{আঁকা (ā̃kā)}]{\bn{আঁকা} (ā̃kā) — to draw}

\label{lesson:BN-C106-aka}



\begin{quote}
Before the new one: say the Bengali for to throw, then the Bengali for to press.
\end{quote}

\subsection*{You'll want to know: \bn{আঁকা}}
\textbf{\bn{আঁকা}} — \emph{ā̃kā} — \textquotedblleft{}to draw\textquotedblright{}.

\begin{cousinweb}
\textbf{\bn{আঁকা}} (\emph{āṁkā}) is a tadbhava, a word inherited from Sanskrit and reshaped by sound change. It comes from \emph{aṅkayati}, \textquotedblleft{}marks, brands\textquotedblright{} (CDIAL 104).
\end{cousinweb}

\begin{grammarlens}[title={What you've built}]
One more word for this chapter.
\end{grammarlens}

\subsection*{Your turn}
\begin{itemize}
\raggedright
  \item \emph{Say it:} \emph{ā̃kā}
  \item \emph{Say it:} \emph{ā̃kā}, once more
  \item \emph{Say it:} say \emph{ā̃kā}
  \item \emph{Recall:} say the Bengali for to throw, then the Bengali for to press, then say \emph{ā̃kā} again
\end{itemize}

\begin{tcolorbox}[breakable,colback=teal!4,colframe=teal!35!black,title={Before you move on}]
What does \bn{আঁকা} mean? (\textquotedblleft{}To draw\textquotedblright{}.) Say it once more.
\end{tcolorbox}
\section[\texorpdfstring{\bn{ফেরা} (pherā)}{ফেরা (pherā)}]{\bn{ফেরা} (pherā) — to return}

\label{lesson:BN-C106-phera}



\begin{quote}
Before the new one: say the Bengali for to press, then the Bengali for to draw.
\end{quote}

\subsection*{You'll want to know: \bn{ফেরা}}
\textbf{\bn{ফেরা}} — \emph{pherā} — \textquotedblleft{}to return\textquotedblright{}.

\begin{cousinweb}
\textbf{\bn{ফেরা}} (\emph{pherā}) is a tadbhava, a word inherited and reshaped by sound change. It comes from \emph{*pherayati}, \textquotedblleft{}makes move, turns\textquotedblright{}, a form CDIAL reconstructs from its descendants (CDIAL 9078). CDIAL lists Bengali \emph{pherā}, \textquotedblleft{}to turn\textquotedblright{}.
\end{cousinweb}

\begin{grammarlens}[title={What you've built}]
One more word for this chapter.
\end{grammarlens}

\subsection*{Your turn}
\begin{itemize}
\raggedright
  \item \emph{Say it:} \emph{pherā}
  \item \emph{Say it:} \emph{pherā}, once more
  \item \emph{Say it:} say \emph{pherā}
  \item \emph{Recall:} say the Bengali for to press, then the Bengali for to draw, then say \emph{pherā} again
\end{itemize}

\begin{tcolorbox}[breakable,colback=teal!4,colframe=teal!35!black,title={Before you move on}]
What does \bn{ফেরা} mean? (\textquotedblleft{}To return\textquotedblright{}.) Say it once more.
\end{tcolorbox}
\section[\texorpdfstring{\bn{চালানো} (chālāno)}{চালানো (chālāno)}]{\bn{চালানো} (chālāno) — to drive, to run (a machine)}

\label{lesson:BN-C106-chalano}



\begin{quote}
Before the new one: say the Bengali for to draw, then the Bengali for to return.
\end{quote}

\subsection*{You'll want to know: \bn{চালানো}}
\textbf{\bn{চালানো}} — \emph{chālāno} — \textquotedblleft{}to drive, to run (a machine)\textquotedblright{}.

\begin{cousinweb}
\textbf{\bn{চালানো}} (\emph{chālāno}) is a tadbhava, a word inherited and reshaped by sound change. It comes from \emph{*calyati}, \textquotedblleft{}goes\textquotedblright{}, a form CDIAL reconstructs from its descendants (CDIAL 4721). \emph{Chālāno} is the causing form of a verb CDIAL lists as Bengali \emph{cālā}, \textquotedblleft{}to be moved\textquotedblright{}.
\end{cousinweb}

\begin{grammarlens}[title={What you've built}]
One more word for this chapter.
\end{grammarlens}

\subsection*{Your turn}
\begin{itemize}
\raggedright
  \item \emph{Say it:} \emph{chālāno}
  \item \emph{Say it:} \emph{chālāno}, once more
  \item \emph{Say it:} say \emph{chālāno}
  \item \emph{Recall:} say the Bengali for to draw, then the Bengali for to return, then say \emph{chālāno} again
\end{itemize}

\begin{tcolorbox}[breakable,colback=teal!4,colframe=teal!35!black,title={Before you move on}]
What does \bn{চালানো} mean? (\textquotedblleft{}To drive, to run (a machine)\textquotedblright{}.) Say it once more.
\end{tcolorbox}
